\section{贪污贿赂罪 - 受贿罪、行贿罪}
\begin{enumerate}
  \item \zd{先找出行贿人的行贿意图}（到底想找谁办事）。\al{甲的丈夫入狱，乙对甲说给我 50W 全部用来帮你打点。实则只想用 10W。甲给了乙 50W，并称如果丈夫出来了，另有重谢。乙 10W 给了丙让丙送给丁，丁当场拒收}。
        \begin{enumerate}
          \item 行贿，标的是权钱交易，买权。\al{甲想找的其实是丁}。
          \item 斡旋，标的是斡旋行为。\al{甲让国乙帮忙捞人 \dui}。
          \item 对有影响力的人行贿，需要答应。\al{甲让乙帮忙捞人 \dui}。
        \end{enumerate}
  \item \kd{受贿}。还要重要，罪名群。
        \begin{enumerate}
          \item 实行行为。交易关系。\zd{权钱交易、不正当的对价关系}。
          \item 交易标的。财物、权（给人办事）。\al{帮官员甲付清装修款 50W、嫖资 \dui。动用关系帮甲升任、股权 \cuo}。
          \item 交易方式。
                \begin{enumerate}
                  \item 事前受贿。有了交易约定。
                        \begin{enumerate}
                          \item （主动）索贿，成立 \yjt 实施索贿行为（\zd{会包含敲诈勒索}）。既遂 \yjt 把钱拿到手。
                          \item （被动）收受。成立 \yjt \zd{许诺办事}。既遂 \yjt（自己 \hz 第三人）收到钱，\zd{事办的如何不问}。
                          \item \zd{给出“许诺”即可}。\al{甲虚假承诺利用职务便利为乙牟利，收取乙财物 \dui}。
                        \end{enumerate}
                  \item 事后受贿。正常办事，没想收钱。不要盯着办事的时候有没有收钱故意。\al{雷政富}。
                        \begin{enumerate}
                          \item \al{乙退休前利用职务便利为甲谋取了不正当利益，退休后收受了甲 10W \dui}。
                          \item \al{离职}。
                        \end{enumerate}
                \end{enumerate}
          \item 变相受贿。\al{甲明知近亲属利用自己的职务行为受贿 \dui。不明知 \yjt 家属利用影响力受贿}。
          \item 既遂标准。\zd{接受他人财物，建立自己的占有} \yjt 不要求取得所有权，后续能处分、支配即可。
                \begin{enumerate}
                  \item \al{甲到局长乙做客，塞银行卡到沙发下 \dui}。
                  \item \al{甲给局长乙一张银行卡，并告知密码，直至被抓乙一直未动里面的钱 \dui}。
                  \item \al{钱先放你那儿吧 \dui}。
                  \item \al{甲收到银行卡，卡内有资金，未设密码 \dui}。
                  \item \al{甲收到银行卡，卡内有资金，甲不知道密码，无法支配卡内资金 \cuo}。
                  \item \al{甲收到银行卡，卡内有资金，甲知道密码，但存的定期 \dui}。
                  \item \al{甲收到银行卡，卡里没钱 \cuo}。
                \end{enumerate}
          \item 定罪数额。\zd{要主客观一致}。详细见财产犯罪 - 既遂和数额部分。
                \begin{enumerate}
                  \item \zd{3W} \yjt 数额较大。
                  \item 20W \yjt 数额巨大，加重处罚。
                  \item 300W \yjt 数额特别巨大，加重处罚。
                  \item \al{乙给甲送茶叶，里面塞了 3W 现金，甲\xhx{不知情}。不构成}。
                  \item \al{乙送甲银行卡，甲问多少钱，乙担心甲不收，\xhx{谎称} 3W，于是甲收下，实际 20W。3W 既遂}。
                  \item \al{乙答应送甲 3W，结果甲打开只收到了 1000。3W 未遂}。
                  \item \al{乙答应送甲 100W，结果甲打开只收到了 3W。3W 既遂，100W 未遂，想象竞合择一重}。
                  \item \al{甲花 5000 购买仿制古董赠与乙，乙以为是价值 20W 的真品。甲不构成行贿，乙 20W 受贿未遂}。
                  \item \al{甲向乙行贿 5W，当晚被丙偷走，后查明 5W 全为假币。甲未达行贿起点、乙 5W 未遂、丙盗窃}。
                  \item \al{100W 的银行卡，案发时，卡内本息 150W。100W 既遂}。
                  \item \al{签发 999W 上限的支票。999W 既遂}。
                \end{enumerate}
          \item 罪数。
                \begin{enumerate}
                  \item 收钱办事构成渎职罪，\zd{原则数罪并罚}。
                  \item 例外。\zd{司法人员执行渎职}（\al{徇私枉法、民事行政枉法裁判、执行判决裁定失职、执行判决裁定滥用职权。择一重）}。
                  \item vs 贪污。\zd{钱财来源属性}。公家 \yjt 贪污。私人 \yjt 受贿。\zd{既受贿又贪污，并罚}。
                        \begin{enumerate}
                          \item \al{甲幼儿园符合 10W 补助条件，局长乙称可以给批，但要分其 4W}
                          \item \al{甲幼儿园符合 10W 补助条件，局长乙称可以给批 14W，但要分其 4W}。
                          \item \al{甲幼儿园不符合 10W 补助条件，送局长乙 4W}。
                          \item 减少债务行为。\al{交警对超载司机说只要交 1000 元罚款的一半，就放你出去。
                                  交警受贿（500）、贪污（1000），司机行贿（500）。
                                  只存在一笔账，羊毛出在羊身上，重吸收轻}。
                        \end{enumerate}
                  \item vs 敲诈勒索、诈骗。中立关系，可以想象竞合。
                \end{enumerate}
        \end{enumerate}
  \item \kd{行贿}。\zd{为了谋取不正当利益}。\al{国乙公司欠甲 100W 工程款，甲行贿丙尽快发钱 \cuo}。
        \begin{enumerate}
          \item 行为类型。
                \begin{enumerate}
                  \item 主动给钱。\zd{钱送出去既遂}，不要求实际获得不正当利益。
                  \item 被动给钱/被勒索给钱。\zd{\blx{实际谋取到不正当利益则构成行贿}}。\al{你不给我 100W，你项目就卡在这}。
                \end{enumerate}
          \item 不正当利益范围扩大解释。包括非法利益、不公平的竞争优势、为第三人谋取等。
                \begin{enumerate}
                  \item \al{考研复试送钱}。
                  \item \al{甲为了是国乙尽快支付工程款，对其送钱 \cuo \yjt 目的正当}。
                  \item \al{甲有求于国乙，明知其妻公司严重亏损，没有扭亏为盈的可能，仍投资了 500W，不久该公司破产}。
                \end{enumerate}
          \item 数额。等号原理。行贿数额 = 受贿数额。基本同受贿。
        \end{enumerate}
  \item 单位问题。个对单、单对个、单对单行贿/受贿，可排列组合。
  \item \zd{所有的行贿、斡旋受贿都要求谋取不正当利益，其他（受贿）不做要求}。
\end{enumerate}
